% TOP_PROBLEM: {"id": 3141, "title": "3-Edge-Coloring Conjecture", "category_id": 3, "category": {"id": 3, "name": "graph_theory", "display_name": "Graph Theory", "description": "Problems involving graphs, networks, and their properties.", "slug": "graph-theory", "order_index": 3, "created_at": "2026-07-31T15:26:25.671Z"}, "status": "open", "problem_number": "OPG-60046", "set_id": 12, "statusQualification": "Uses simple starting cubic graphs and retains parallel edges after suppression; the equivalent deleted-graph flow condition makes the reduction precise."}
% FIRST_POSTED: 2026-10-01
% ENTRY_KIND: statement_only
% UnsolvedMath ID 3141; source catalogue code OPG-60046
% !TeX program = lualatex
% Definition and sources only; this file is not an LLM solution attempt.
\documentclass[11pt,a4paper]{article}
\usepackage[margin=25mm]{geometry}
\usepackage{fontspec}
\setmainfont{lmroman10-regular.otf}[BoldFont=lmroman10-bold.otf,
 ItalicFont=lmroman10-italic.otf,BoldItalicFont=lmroman10-bolditalic.otf]
\usepackage{amsmath,amssymb,mathtools}
\usepackage{unicode-math}
\setmathfont{latinmodern-math.otf}
\AtBeginDocument{\let\setminus\smallsetminus}
\usepackage{xurl,hyperref}
\hypersetup{colorlinks=true,urlcolor=blue}
\setcounter{secnumdepth}{0}
\setlength{\parindent}{0pt}
\setlength{\parskip}{5pt}
\setlength{\emergencystretch}{3em}
\begin{document}
\section*{3-Edge-Coloring Conjecture}\label{3141}
{\small UnsolvedMath ID 3141 · OPG-60046 · Graph Theory · OpenGarden}
\subsection{Definitions and mathematical statement}
Let $G$ be a finite connected cubic graph with more than two vertices. Here cubic means that every vertex has degree three, and the starting graph is simple. A proper three-edge-coloring assigns one of three colors to each edge so that edges sharing a vertex receive distinct colors; consequently every vertex sees all three colors.

For an edge $e=uv$, delete $e$. The vertices $u$ and $v$ now have degree two. Suppress each of them: replace the two remaining incident edges through that vertex by one edge joining their other ends. Denote the resulting cubic multigraph by $G\mathbin{\ominus}e$. Parallel edges created by suppression are retained as distinct edges. A loop, if created in a more general multigraph version, cannot receive a proper edge color. The question asks whether every three-edge-colorable $G$ has at least one edge $e$ for which $G\mathbin{\ominus}e$ is again three-edge-colorable.

Equivalently, the deleted graph $G-e$ should admit a nowhere-zero flow over $\mathbb Z_2\times\mathbb Z_2$: assign a nonzero group element to every remaining edge so that the sum of incident values at each vertex is zero. At a degree-two vertex the two values must agree, making suppression legitimate; at a degree-three vertex the three nonzero values must all differ. The coloring after deletion and suppression is allowed to be entirely different from any initially supplied coloring of $G$.
\subsection{Short English statement}
Does every connected three-edge-colorable cubic graph have an edge whose deletion and suppression leaves another three-edge-colorable cubic multigraph?
\subsection{Sources}
\begin{itemize}
\item[S1] ULAM.AI and the UnsolvedMath Contributors, supplied problems.json, record 3141 (OPG-60046), OpenGarden. Catalogue identity and source transcription; CC BY 4.0.
\url{https://huggingface.co/datasets/ulamai/UnsolvedMath}.
\item[S2] Open Problem Garden, 3-Edge-Coloring Conjecture. Original problem and discussion.
\url{http://www.openproblemgarden.org/op/3_edge_coloring_conjecture}.
\end{itemize}
\end{document}
